(a) Se a taxa de contagem é $c$ pessoas por unidade de tempo, o tempo necessário para contar a população $P(t)$ é
\[
\tau = \frac{P(t)}{c}.
\]
(b) O aumento da população durante esse tempo é aproximadamente $\lambda\tau P(t)$, de modo que o tempo para contar o aumento é
\[
\tau_{\text{aumento}} \approx \frac{\lambda P(t)}{c} \tau = \frac{\lambda P(t)^{2}}{c^{2}},
\]
mas, na prática, o tempo adicional é do mesmo tipo e depende da taxa de contagem e da taxa de crescimento. 